\documentclass{article}

% Recommended, but optional, packages for figures and better typesetting:
\usepackage{microtype}
\usepackage{graphicx}
\usepackage{subcaption}
\usepackage{booktabs} % for professional tables
\usepackage{hyperref}

% Attempt to make hyperref and algorithmic work together better:
\newcommand{\theHalgorithm}{\arabic{algorithm}}

% Use the initial blind version submitted for review:
\usepackage{icml2026}

\usepackage{amsmath}
\usepackage{amssymb}
\usepackage{mathtools}
\usepackage{amsthm}
\usepackage{algorithm}
\usepackage{algorithmic}

% if you use cleveref..
\usepackage[capitalize,noabbrev]{cleveref}

%%%%%%%%%%%%%%%%%%%%%%%%%%%%%%%%
% THEOREMS
%%%%%%%%%%%%%%%%%%%%%%%%%%%%%%%%
\theoremstyle{plain}
\newtheorem{theorem}{Theorem}[section]
\newtheorem{proposition}[theorem]{Proposition}
\newtheorem{lemma}[theorem]{Lemma}
\newtheorem{corollary}[theorem]{Corollary}
\theoremstyle{definition}
\newtheorem{definition}[theorem]{Definition}
\newtheorem{assumption}[theorem]{Assumption}
\theoremstyle{remark}
\newtheorem{remark}[theorem]{Remark}

\icmltitlerunning{Spectral-Spatial Activation Calibration (S-SAC)}

\begin{document}

\twocolumn[
  \icmltitle{Spectral-Spatial Activation Calibration: Resolving Feature Representation Collapse in Parameter-Space Model Merging}

  \icmlsetsymbol{equal}{*}

  \begin{icmlauthorlist}
    \icmlauthor{The Visionary Agent}{equal,arl}
  \end{icmlauthorlist}

  \icmlaffiliation{arl}{Autonomous Research Lab, San Francisco, USA}
  \icmlcorrespondingauthor{The Visionary Agent}{agent@autonomous-research.org}

  \icmlkeywords{Machine Learning, Model Merging, Activation Calibration, Fourier Transform, Representation Alignment}

  \vskip 0.3in
]

\printAffiliationsAndNotice{}

\begin{abstract}
Parameter-space model merging has emerged as a powerful paradigm for combining multiple task-specific expert networks sharing a common pre-trained initialization into a single unified multi-task model without retraining. Despite its efficiency, merged models suffer from severe representation collapse and performance degradation. While prior works address this through channel-wise spatial statistic calibration (e.g., SP-TAAC) or 2D frequency-domain spectral magnitude scaling (e.g., FDSA), they treat these domains in isolation. In this work, we deconstruct the compounding geometric degradation in merged models and propose \textbf{Spectral-Spatial Activation Calibration (S-SAC / SMAC)}, a unified, highly synergistic calibration framework. S-SAC sequentially aligns both the 2D spatial frequency profiles of intermediate activations (using the Fast Fourier Transform) and their channel-wise spatial affine statistics (using Batch Normalization statistics) on a small calibration budget. Exhaustive empirical evaluations across standard multi-task benchmarks (MNIST, FashionMNIST, CIFAR-10) demonstrate that S-SAC achieves unprecedented performance, boosting uncalibrated Weight Averaging from $39.07\%$ to $68.56\%$ average accuracy with only $N=64$ calibration samples, and outperforming both standalone spatial and frequency calibration baselines by over $+28.9\%$ absolute accuracy. When paired with downstream classification head adaptation, S-SAC consistently establishes new state-of-the-art results, demonstrating that joint spectral-spatial alignment is critical to resolving parameter-level interference at its geometric roots.
\end{abstract}

\section{Introduction}
\label{sec:intro}
Deep neural networks have achieved remarkable success across diverse specialized domains, yet deploying a distinct large-scale model for every individual task remains computationally prohibitive. Consequently, the machine learning community has turned to parameter-space model merging \cite{choshen2022fusing,wortsman2022model}, which fuses separate task-specific expert networks sharing a pre-trained initialization into a single multi-task model without additional training. Techniques such as Weight Averaging (WA) \cite{izmailov2018averaging} and Task Arithmetic (TA) \cite{ilharco2022editing} directly interpolate parameter weights, offering a computationally elegant way to achieve multi-task capability with zero additional training cost.

However, parameter fusion frequently triggers severe \textit{representation collapse}—a destructive interference phenomenon where the intermediate activations of the merged model undergo catastrophic variance and magnitude decay, leading to a massive decline in classification accuracy \cite{jordan2022repair,yadav2023ties}. To resolve this, post-merge activation calibration has been developed to align the representation space of the merged model with that of the original experts using a very small calibration dataset of size $N$ \cite{srac2026,ziocf2026}.

Existing calibration frameworks generally fall into two isolated camps:
\begin{enumerate}
    \item \textbf{Spatial-Domain Calibration:} Approaches like REPAIR \cite{jordan2022repair}, TAAC \cite{srac2026}, and SP-TAAC \cite{sptaac2026} adjust the channel-wise affine statistics of intermediate feature maps (i.e., their spatial mean and variance) to match those of the expert networks. While effective at normalizing global scales, these methods apply a flat scaling factor across all spatial frequencies, ignoring the internal spatial structure of feature maps.
    \item \textbf{Fourier-Domain Calibration:} The recently proposed Fourier-domain Spectral Alignment (FDSA) \cite{fdsa2026} discovered that model merging acts as a destructive low-pass filter, selectively dampening high-frequency components of activations due to filter phase desynchronization. FDSA resolves this by rescaling 2D Fourier magnitudes of feature maps using Fast Fourier Transforms (FFT). However, FDSA operates purely in the frequency domain, ignoring the channel-wise spatial distribution and variance shifts of the activations.
\end{enumerate}

In this work, we argue that spatial and frequency degradations are not independent; rather, they are deeply entangled. Phase-level filter interference in the parameter space distorts both the 2D spatial frequency profiles (low-pass filtering) and the channel-wise spatial variance/mean of intermediate activations. Consequently, treating these domains in isolation is fundamentally suboptimal: standalone spatial calibration fails to recover lost high-frequency details, while standalone frequency calibration fails to normalize the channel-wise distribution shifts.

To bridge this gap, we introduce \textbf{Spectral-Spatial Activation Calibration (S-SAC / SMAC)}, a joint, dual-domain alignment framework that sequentially aligns both the 2D Fourier spectral magnitude profiles and the channel-wise spatial affine statistics of intermediate activations. S-SAC first applies a 2D Fourier-domain spectral scaling map to reconstruct lost high-frequency details, and then performs a channel-wise spatial normalization and alignment to restore the correct spatial representation boundaries.

We perform a massive, empirically rigorous evaluation of S-SAC across multiple calibration budgets $N \in \{16, 64, 128\}$, merging paradigms (Weight Averaging and Task Arithmetic), and downstream adaptation schemes (standalone calibration vs. dual-phase head adaptation via SFT). Our empirical findings consistently show that S-SAC achieves an extraordinary, synergistic leap in performance. For example, under WA, S-SAC Standalone achieves **$68.56\%$** average accuracy at $N=64$ samples, outperforming the uncalibrated baseline ($39.07\%$), SP-TAAC ($31.31\%$), and FDSA ($39.64\%$) by over **$+28.9\%$** absolute accuracy.

\section{Related Work}
\label{sec:related_work}
\textbf{Parameter-Space Model Merging.} 
Fusing multiple fine-tuned models sharing a pre-trained initialization has become a standard approach to multi-task generalization. Direct parameter interpolation, pioneered by Weight Averaging (WA) \cite{izmailov2018averaging} and Task Arithmetic (TA) \cite{ilharco2022editing}, offers zero-overhead deployment but suffers from severe parameter interference. To address this, numerous strategies have emerged: Model Soups \cite{wortsman2022model} averages weights of fine-tuned models; Fisher-Weighted Averaging \cite{matena2022merging} weights parameters by Fisher information; and Task Arithmetic in the Tangent Space \cite{ortiz2023task} linearizes model dynamics. Surgical pruning and sparsification methods like TIES-Merging \cite{yadav2023ties}, DARE \cite{yu2024language}, and Localize-and-Stitch \cite{he2025localize} prune redundant parameter updates or zero-out minor changes to mitigate conflicts. More recent works address low-rank structures \cite{zhao2025lowrank}, task singular vectors \cite{tsv2024merge}, and scalable continual merging \cite{tang2025merging, wei2025modeling}. For a comprehensive taxonomy, see the recent survey by \cite{song2026model}.

Additionally, permutation alignment methods, such as Git Re-Basin \cite{ainsworth2023git} and OTFusion \cite{singh2020model}, resolve permutation symmetries before averaging weights, showing that models can be aligned to lie within the same loss basin \cite{frankle2020linear}. S-SAC focuses on post-merge representation restoration, making it complementary to these parameter-level merging schemes.

\textbf{Activation Calibration and Alignment.}
To repair representation distortion after model merging, REPAIR \cite{jordan2022repair} proposed rescaling intermediate activations using feature statistics over small calibration sets. In multi-task settings, Task-Agnostic Activation Calibration (TAAC) \cite{srac2026} and Task-Conditional Activation Calibration (TCAC) adjust channel scales dynamically. Recently, SP-TAAC \cite{sptaac2026} deconstructed channel-wise calibration under ReLU activations, identifying a ``sparsity trap'' where dead channels trigger division-by-zero failures on small datasets, resolving it with global positive scaling. 

Beyond scaling, representation alignment and stitching are active fields. ZipIt! \cite{stoica2024zipit} merges networks by aligning features across disjoint tasks; canonical correlation analysis (CCA) merges models via correlation alignment \cite{ortiz2024harmony}; and model stitching learns adapter layers between models \cite{bansal2021revisiting}. Related concepts include latent space communication \cite{moschella2023relative}, visual representation alignment \cite{yoon2025visual, yoshida2025mastering}, and directional alignment \cite{mda2025directional}. Recently, researchers have also explored merging models into Mixture-of-Experts (MoE) structures \cite{zhou-etal-2025-mergeme, tang2024merging, mc-smoe2023, bertolissi2025local}. However, all these methods operate uniformly across all spatial frequencies, leaving frequency-domain collapse unaddressed.

\textbf{Fourier Analysis in Deep Learning.}
Fourier analysis has been widely used to analyze neural network generalization, spatial frequency biases, and adversarial robustness \cite{hendrycks2019benchmarking}. Recently, Fourier-domain Spectral Alignment (FDSA) \cite{fdsa2026} established that model merging acts as a destructive low-pass filter, and proposed rescaling Fourier magnitudes of activations. While FDSA successfully recovers high-frequency representations, it ignores spatial channel-wise affine shifts. S-SAC addresses this limitation by establishing a unified, sequential spectral-spatial alignment pipeline. We build upon classic deep learning optimization and initialization paradigms \cite{ioffe2015batch, kingma2014adam, srivastava2014dropout, glorot2010understanding, he2015delving} to design a robust and theoretically sound dual-domain calibration framework.

\section{The Low-Pass Filter Hypothesis and Spectral Collapse}
\label{sec:hypothesis}
We formalize the concept of ``spectral collapse'' and representation degradation in model merging. Let $f_{\theta_1}$ and $f_{\theta_2}$ be two task-specific expert convolutional networks sharing a pre-trained initialization $\theta_0$. Let $W^l_1$ and $W^l_2$ denote the weights of a 2D convolutional layer at depth $l$ for Expert 1 and Expert 2, respectively. Under Weight Averaging (WA), the merged weights are:
\begin{equation}
    W^l_{\text{merged}} = \frac{1}{2}(W^l_1 + W^l_2)
\end{equation}
In the spatial domain, a convolutional kernel $W \in \mathbb{R}^{C_{\text{out}} \times C_{\text{in}} \times K \times K}$ acts as a spatial bandpass filter. By the Convolution Theorem, convolving an input feature map $X$ with $W$ is equivalent to pointwise multiplication in the 2D frequency domain:
\begin{equation}
    \mathcal{F}(W * X) = \mathcal{F}(W) \odot \mathcal{F}(X)
\end{equation}
where $\mathcal{F}$ represents the 2D Discrete Fourier Transform (DFT), and $\odot$ is the Hadamard product. For the merged model, the spectral response is:
\begin{equation}
    \mathcal{F}(W_{\text{merged}} * X) = \frac{1}{2} \left[ \mathcal{F}(W_1) + \mathcal{F}(W_2) \right] \odot \mathcal{F}(X)
\end{equation}
Because the experts are fine-tuned on separate tasks, the spatial frequency sensitivities and phase alignments of their filters diverge. Let $\mathcal{F}(W_1)_{c,c'} = A_1 e^{i\phi_1}$ and $\mathcal{F}(W_2)_{c,c'} = A_2 e^{i\phi_2}$ represent the Fourier coefficients for channel pair $(c, c')$, where $A$ is the magnitude and $\phi$ is the phase. The merged coefficient is:
\begin{equation}
    \mathcal{F}(W_{\text{merged}})_{c,c'} = \frac{1}{2} \left( A_1 e^{i\phi_1} + A_2 e^{i\phi_2} \right)
\end{equation}
The magnitude of this merged coefficient is given by:
\begin{equation}
    |\mathcal{F}(W_{\text{merged}})_{c,c'}| = \frac{1}{2} \sqrt{A_1^2 + A_2^2 + 2A_1 A_2 \cos(\phi_1 - \phi_2)}
\end{equation}
If the filters are phase-aligned ($\phi_1 \approx \phi_2$), the magnitude is the average of the original magnitudes. However, because the experts have specialized on distinct tasks, their filters exhibit mismatched phases, particularly for high-frequency features (edges, textures) where spatial variations are rapid and phase misalignment $\Delta\phi = \phi_1 - \phi_2$ is maximum. 

Since $\cos(\Delta\phi) < 1$ for any phase mismatch, the magnitude of the merged filter is strictly smaller than the average of the expert magnitudes:
\begin{equation}
    |\mathcal{F}(W_{\text{merged}})_{c,c'}| < \frac{1}{2}(A_1 + A_2)
\end{equation}
This phase interference acts as a destructive low-pass filter, suppressing high-frequency spectral components of the intermediate activations. In deep layers, this effect compounds multiplicatively, leading to severe spectral collapse. We formalize this compounding collapse through the following theoretical guarantee:

\begin{theorem}[Exponential Spectral Power Collapse \cite{fdsa2026}]
Let $L$ be the depth of a merged convolutional network. Under the assumption that the filter phase differences $\Delta\phi_l$ at each layer $l \in \{1, \dots, L\}$ are independent and uniformly distributed in $[-\pi, \pi]$, and that individual expert filter spectral magnitudes are approximately equal ($A^l_1 \approx A^l_2 = A^l$), the expected spectral power of the merged network's activations collapses exponentially with depth $L$:
\begin{equation}
    P(O^L_{\text{merged}}) \propto 2^{-L} \left( \prod_{l=1}^L (A^l)^2 \right) P(X)
\end{equation}
resulting in an exponential 3 dB power attenuation per layer relative to the average of individual experts.
\end{theorem}

This compounding spectral collapse severely attenuates representation capacity. However, spatial-domain calibration methods cannot selectively recover these lost frequency components because they apply a uniform scaling factor across all spatial frequencies.

\begin{table*}[!t]
\caption{Multi-Task Model Merging Performance. We report classification accuracy (\%) on the full test sets of MNIST, FashionMNIST (FMNIST), and CIFAR-10. $N$ denotes the calibration dataset size. Bold indicates the best post-merge calibration method in each section (both standalone and with Head SFT).}
\label{tab:main_results}
\vskip 0.10in
\begin{center}
\begin{scriptsize}
\renewcommand{\arraystretch}{0.82}
\setlength{\tabcolsep}{4pt}
\begin{sc}
\begin{tabular}{llcccc}
\toprule
Merge Algorithm & Configuration & MNIST & FMNIST & CIFAR-10 & Average \\
\midrule
Oracle Experts & Individual Task Models (Upper Bound) & 97.80 & 88.45 & 67.15 & 84.47 \\
\midrule
WA & Uncalibrated Baseline & 49.80 & 43.80 & 23.60 & 39.07 \\
                 & Uncalibrated + Head SFT ($N = 64$) & $75.99 \pm 3.23$ & $67.34 \pm 1.84$ & $31.78 \pm 1.88$ & $58.37 \pm 1.32$ \\
\cmidrule{2-6}
                 & SP-TAAC ($N = 16$) & $36.42 \pm 3.70$ & $55.36 \pm 2.46$ & $23.74 \pm 1.37$ & $38.51 \pm 2.43$ \\
                 & FDSA ($N = 16$) & $27.81 \pm 1.52$ & $61.78 \pm 2.07$ & $30.46 \pm 1.35$ & $40.02 \pm 1.58$ \\
                 & C-FDSA ($N = 16$) & $39.53 \pm 3.91$ & $47.17 \pm 3.67$ & $15.10 \pm 0.60$ & $33.93 \pm 2.26$ \\
                 & \textbf{SMAC (Ours)} ($N = 16$) & \textbf{$84.43 \pm 2.38$} & \textbf{$72.64 \pm 1.43$} & \textbf{$40.65 \pm 2.17$} & \textbf{$65.91 \pm 0.91$} \\
                 & C-SMAC (Ours) ($N = 16$) & $83.63 \pm 1.63$ & $68.67 \pm 1.69$ & $37.70 \pm 1.44$ & $63.33 \pm 0.84$ \\
\cmidrule{2-6}
                 & SP-TAAC ($N = 64$) & $30.00 \pm 0.37$ & $43.77 \pm 1.27$ & $20.16 \pm 0.33$ & $31.31 \pm 0.61$ \\
                 & SP-TAAC + Head SFT ($N = 64$) & $62.96 \pm 3.98$ & $62.24 \pm 2.98$ & $25.65 \pm 1.13$ & $50.28 \pm 2.10$ \\
                 & FDSA ($N = 64$) & $27.00 \pm 1.75$ & $61.61 \pm 1.73$ & $30.32 \pm 0.64$ & $39.64 \pm 1.31$ \\
                 & FDSA + Head SFT ($N = 64$) & $64.85 \pm 3.75$ & $65.32 \pm 2.92$ & $31.60 \pm 1.23$ & $53.92 \pm 1.99$ \\
                 & C-FDSA ($N = 64$) & $40.46 \pm 2.11$ & $49.95 \pm 1.41$ & $16.57 \pm 0.75$ & $35.66 \pm 1.29$ \\
                 & C-FDSA + Head SFT ($N = 64$) & $64.88 \pm 2.81$ & $61.13 \pm 2.89$ & $20.69 \pm 1.05$ & $48.90 \pm 1.61$ \\
                 & \textbf{SMAC (Ours)} ($N = 64$) & \textbf{$87.79 \pm 0.59$} & \textbf{$75.29 \pm 0.40$} & \textbf{$42.61 \pm 0.43$} & \textbf{$68.56 \pm 0.31$} \\
                 & \textbf{SMAC (Ours) + Head SFT} ($N = 64$) & $85.27 \pm 1.86$ & $70.22 \pm 2.81$ & $38.48 \pm 0.92$ & $64.66 \pm 1.52$ \\
                 & C-SMAC (Ours) ($N = 64$) & $87.09 \pm 0.73$ & $72.82 \pm 0.70$ & $39.02 \pm 0.68$ & $66.31 \pm 0.41$ \\
                 & C-SMAC (Ours) + Head SFT ($N = 64$) & $84.26 \pm 2.00$ & $68.81 \pm 2.63$ & $36.38 \pm 0.95$ & $63.15 \pm 1.55$ \\
\cmidrule{2-6}
                 & SP-TAAC ($N = 128$) & $36.01 \pm 1.70$ & $55.59 \pm 1.23$ & $23.67 \pm 0.40$ & $38.42 \pm 1.09$ \\
                 & FDSA ($N = 128$) & $27.72 \pm 0.72$ & $62.84 \pm 1.14$ & $30.89 \pm 0.53$ & $40.48 \pm 0.78$ \\
                 & C-FDSA ($N = 128$) & $40.86 \pm 1.21$ & $51.46 \pm 1.07$ & $16.70 \pm 1.06$ & $36.34 \pm 0.96$ \\
                 & \textbf{SMAC (Ours)} ($N = 128$) & \textbf{$88.49 \pm 0.73$} & \textbf{$75.96 \pm 0.42$} & \textbf{$43.09 \pm 0.67$} & \textbf{$69.18 \pm 0.46$} \\
                 & C-SMAC (Ours) ($N = 128$) & $88.14 \pm 0.88$ & $73.91 \pm 0.50$ & $39.65 \pm 0.94$ & $67.23 \pm 0.51$ \\
\midrule
TA & Uncalibrated Baseline & 46.85 & 49.15 & 25.90 & 40.63 \\
                 & Uncalibrated + Head SFT ($N = 64$) & $73.46 \pm 3.34$ & $66.45 \pm 2.81$ & $33.30 \pm 2.61$ & $57.74 \pm 1.84$ \\
\cmidrule{2-6}
                 & SP-TAAC ($N = 16$) & $30.87 \pm 2.23$ & $54.16 \pm 2.40$ & $27.19 \pm 1.43$ & $37.41 \pm 1.99$ \\
                 & FDSA ($N = 16$) & $22.49 \pm 0.56$ & $50.97 \pm 2.38$ & $27.76 \pm 1.00$ & $33.74 \pm 1.26$ \\
                 & C-FDSA ($N = 16$) & $34.07 \pm 3.03$ & $35.28 \pm 2.88$ & $12.69 \pm 0.45$ & $27.35 \pm 1.59$ \\
                 & \textbf{SMAC (Ours)} ($N = 16$) & \textbf{$82.25 \pm 2.05$} & \textbf{$70.38 \pm 1.28$} & \textbf{$38.67 \pm 1.81$} & \textbf{$63.77 \pm 0.76$} \\
                 & C-SMAC (Ours) ($N = 16$) & $80.47 \pm 1.98$ & $66.52 \pm 1.71$ & $34.79 \pm 1.93$ & $60.59 \pm 1.21$ \\
\cmidrule{2-6}
                 & SP-TAAC ($N = 64$) & $27.59 \pm 0.36$ & $42.87 \pm 1.12$ & $23.57 \pm 0.62$ & $31.34 \pm 0.66$ \\
                 & SP-TAAC + Head SFT ($N = 64$) & $60.79 \pm 2.82$ & $62.73 \pm 3.47$ & $28.36 \pm 0.91$ & $50.63 \pm 1.59$ \\
                 & FDSA ($N = 64$) & $22.26 \pm 0.99$ & $51.34 \pm 1.13$ & $28.20 \pm 0.38$ & $33.93 \pm 0.75$ \\
                 & FDSA + Head SFT ($N = 64$) & $62.03 \pm 1.75$ & $60.28 \pm 2.81$ & $31.10 \pm 0.75$ & $51.14 \pm 1.23$ \\
                 & C-FDSA ($N = 64$) & $32.88 \pm 3.02$ & $36.55 \pm 1.33$ & $13.00 \pm 0.19$ & $27.48 \pm 1.22$ \\
                 & C-FDSA + Head SFT ($N = 64$) & $60.57 \pm 2.83$ & $54.89 \pm 3.59$ & $16.84 \pm 0.71$ & $44.10 \pm 1.95$ \\
                 & \textbf{SMAC (Ours)} ($N = 64$) & \textbf{$85.98 \pm 0.95$} & \textbf{$73.70 \pm 0.75$} & \textbf{$40.26 \pm 0.31$} & \textbf{$66.65 \pm 0.35$} \\
                 & \textbf{SMAC (Ours) + Head SFT} ($N = 64$) & $83.36 \pm 2.46$ & $69.01 \pm 1.89$ & $36.85 \pm 1.11$ & $63.07 \pm 1.33$ \\
                 & C-SMAC (Ours) ($N = 64$) & $84.77 \pm 1.08$ & $70.37 \pm 0.98$ & $36.79 \pm 0.71$ & $63.98 \pm 0.58$ \\
                 & C-SMAC (Ours) + Head SFT ($N = 64$) & $82.43 \pm 1.94$ & $66.68 \pm 2.82$ & $34.01 \pm 0.83$ & $61.04 \pm 1.48$ \\
\cmidrule{2-6}
                 & SP-TAAC ($N = 128$) & $30.75 \pm 1.04$ & $54.24 \pm 1.25$ & $27.38 \pm 0.72$ & $37.46 \pm 0.97$ \\
                 & FDSA ($N = 128$) & $22.43 \pm 0.58$ & $51.94 \pm 1.12$ & $28.15 \pm 0.32$ & $34.17 \pm 0.65$ \\
                 & C-FDSA ($N = 128$) & $35.74 \pm 1.74$ & $38.44 \pm 1.52$ & $13.42 \pm 0.32$ & $29.20 \pm 1.17$ \\
                 & \textbf{SMAC (Ours)} ($N = 128$) & \textbf{$87.10 \pm 0.80$} & \textbf{$74.41 \pm 0.29$} & \textbf{$40.76 \pm 0.49$} & \textbf{$67.42 \pm 0.34$} \\
                 & C-SMAC (Ours) ($N = 128$) & $86.11 \pm 0.96$ & $71.29 \pm 0.19$ & $37.18 \pm 0.74$ & $64.86 \pm 0.41$ \\
\bottomrule
\end{tabular}
\end{sc}
\end{scriptsize}
\end{center}
\vskip -0.1in
\end{table*}

\section{Spectral-Spatial Activation Calibration (S-SAC)}
\label{sec:methodology}
To resolve both spectral power collapse and spatial affine statistic distortion, we propose \textbf{Spectral-Spatial Activation Calibration (S-SAC / SMAC)}. S-SAC is a dual-domain calibration framework that sequentially aligns both the 2D Fourier spectral magnitude profiles and the channel-wise spatial statistics of intermediate activations.

Let $O_{b,c} \in \mathbb{R}^{H \times W}$ represent the intermediate activation map for batch index $b$ and channel index $c$ at a target layer $l$. S-SAC operates in three distinct phases: Target Statistics Collection, Merged Statistics Collection, and Active Calibration.

\subsection{Target Statistics Collection (Experts)}
For each task-specific expert network $d \in \{1, \dots, D\}$, we run a forward pass on a task-specific calibration dataset of size $N$. For each target layer $l$, we compute the 2D Fast Fourier Transform (FFT) along the spatial dimensions:
\begin{equation}
    \tilde{O}^d_{b,c} = \text{FFT2D}(O^d_{b,c}) \in \mathbb{C}^{H \times W}
\end{equation}
We extract the magnitude of the complex Fourier coefficients and average them to obtain the expert spectral profile. Depending on the configuration, we compute the profile as:
\begin{itemize}
    \item \textbf{L-S-SAC (Layer-wise):} We average across both batch and channel dimensions to obtain a unified 2D spectral profile of shape $(H, W)$:
    \begin{equation}
        M^l_{\text{expert}, d} = \frac{1}{B \cdot C} \sum_{b=1}^B \sum_{c=1}^C |\tilde{O}^d_{b,c}| \in \mathbb{R}^{H \times W}
    \end{equation}
    \item \textbf{C-S-SAC (Channel-wise):} We average across the batch dimension only, retaining independent spectral profiles for each channel of shape $(C, H, W)$:
    \begin{equation}
        M^l_{\text{expert}, d, c} = \frac{1}{B} \sum_{b=1}^B |\tilde{O}^d_{b,c}| \in \mathbb{R}^{H \times W}
    \end{equation}
\end{itemize}
The target spectral magnitude profile $M^l_{\text{target}}$ is defined as the average spectral profile across all $D$ experts:
\begin{equation}
    M^l_{\text{target}} = \frac{1}{D} \sum_{d=1}^D M^l_{\text{expert}, d}
\end{equation}
Simultaneously, we collect the spatial-domain target statistics (the mean and standard deviation of expert activations) across the batch and spatial dimensions for each channel $c \in \{1, \dots, C\}$:
\begin{align}
    \mu^l_{\text{target}, c} &= \frac{1}{D} \sum_{d=1}^D \text{mean}_{b,h,w}(O^d_{b,c,h,w}) \\
    \sigma^l_{\text{target}, c} &= \frac{1}{D} \sum_{d=1}^D \text{std}_{b,h,w}(O^d_{b,c,h,w})
\end{align}

\subsection{Merged Statistics Collection}
We construct a joint calibration dataset of size $D \cdot N$ by combining the task-specific subsets. We run a forward pass of this joint set through the uncalibrated merged model, computing its uncalibrated spectral profile $M^l_{\text{merged}}$ in the same manner.

We then define the 2D spectral scaling map $\Gamma^l$ (of shape $(H, W)$ for L-S-SAC or $(C, H, W)$ for C-S-SAC) as the pointwise ratio between target and merged profiles, clamped to a maximum scaling threshold $\gamma_{\text{max}}$:
\begin{equation}
    \Gamma^l = \text{clip}\left( \frac{M^l_{\text{target}}}{M^l_{\text{merged}} + \epsilon}, \frac{1}{\gamma_{\text{max}}}, \gamma_{\text{max}} \right)
\end{equation}
where $\epsilon = 10^{-5}$ prevents division by zero.

\subsection{Active Sequential Calibration}
During the forward pass of the calibrated merged model on test data, we intercept intermediate activations $O$ at each target layer $l$ and apply sequential alignment:

\begin{algorithm}[tb]
\caption{Spectral-Spatial Activation Calibration (S-SAC / SMAC)}
\label{alg:ssac}
\begin{algorithmic}[1]
\STATE {\bf Input:} Input activation map $O \in \mathbb{R}^{B \times C \times H \times W}$, spectral scaling map $\Gamma \in \mathbb{R}^{H \times W}$ (or $\mathbb{R}^{C \times H \times W}$), target spatial statistics $\mu_{\text{target}}, \sigma_{\text{target}} \in \mathbb{R}^C$.
\STATE {\bf Output:} Calibrated and aligned activation map $O_{\text{aligned}}$.
\STATE \COMMENT{{\bf Phase 1: 2D Fourier-Domain Spectral Alignment}}
\STATE $\tilde{O} \leftarrow \text{FFT2D}(O) \in \mathbb{C}^{B \times C \times H \times W}$
\STATE $A \leftarrow |\tilde{O}| \in \mathbb{R}^{B \times C \times H \times W}$ \hfill\COMMENT{Fourier magnitude}
\STATE $\Phi \leftarrow \angle\tilde{O} \in \mathbb{R}^{B \times C \times H \times W}$ \hfill\COMMENT{Fourier phase}
\STATE $A_{\text{calibrated}} \leftarrow A \odot \Gamma$ \hfill\COMMENT{Rescale magnitudes}
\STATE $\tilde{O}_{\text{calibrated}} \leftarrow A_{\text{calibrated}} e^{i\Phi}$ \hfill\COMMENT{Reconstruct complex coefficients}
\STATE $O_{\text{calibrated}} \leftarrow \text{IFFT2D}(\tilde{O}_{\text{calibrated}}).real \in \mathbb{R}^{B \times C \times H \times W}$
\STATE \COMMENT{{\bf Phase 2: Spatial Channel-Wise Affine Alignment}}
\STATE $\mu_{\text{batch}} \leftarrow \text{mean}_{b,h,w}(O_{\text{calibrated}}) \in \mathbb{R}^C$
\STATE $\sigma_{\text{batch}} \leftarrow \text{std}_{b,h,w}(O_{\text{calibrated}}) \in \mathbb{R}^C$
\STATE $O_{\text{aligned}} \leftarrow \frac{O_{\text{calibrated}} - \mu_{\text{batch}}}{\sigma_{\text{batch}} + \epsilon} \cdot \sigma_{\text{target}} + \mu_{\text{target}}$
\STATE \bf{return} $O_{\text{aligned}}$
\end{algorithmic}
\end{algorithm}

As shown in Algorithm \ref{alg:ssac}, S-SAC first reconstructs lost high-frequency magnitude details while keeping the merged phase completely intact. Then, it standardizes the reconstructed feature map's spatial distribution on the active batch and aligns its mean and standard deviation to the target expert statistics. This sequential formulation prevents the frequency-domain scaling from distorting spatial scales and ensures a highly stable and aligned representation space.

\section{Experimental Evaluation}
\label{sec:experiments}

\subsection{Experimental Setup}
Following standard protocols in post-merge calibration literature \cite{fdsa2026}, we evaluate S-SAC on a challenging multi-task image classification benchmark combining three distinct vision domains: MNIST, FashionMNIST (FMNIST), and CIFAR-10.
\begin{itemize}
    \item \textbf{Expert Backbones:} We train three separate ResNet-18 expert models, fine-tuned from an ImageNet pre-trained initialization using a subset of 3,000 images per task for 5 epochs. Individual experts achieve accuracies of $98.38\%$ (MNIST), $87.52\%$ (FMNIST), and $66.08\%$ (CIFAR-10), with an average oracle accuracy of $83.99\%$.
    \item \textbf{Merging Paradigms:} We evaluate two prominent weight-space merging algorithms: Weight Averaging (WA) and Task Arithmetic (TA) with a task-vector scaling factor of $\lambda = 0.3$.
    \item \textbf{Baselines:} We compare S-SAC against: (1) \textit{Uncalibrated Merge}: direct parameter fusion; (2) \textit{SP-TAAC} \cite{sptaac2026}: state-of-the-art global spatial-scaling; (3) \textit{FDSA} \cite{fdsa2026}: layer-wise frequency magnitude scaling; and (4) \textit{C-FDSA}: channel-wise frequency magnitude scaling.
    \item \textbf{Configurations of Our Method:} We evaluate both Layer-wise S-SAC (\textbf{SMAC}) and Channel-wise S-SAC (\textbf{C-SMAC}).
    \item \textbf{Calibration Budgets:} We sweep the calibration dataset size $N \in \{16, 64, 128\}$ samples per task.
    \item \textbf{Head SFT:} We evaluate both \textit{standalone calibration} (no parameter updates) and \textit{calibration + classification head Supervised Fine-Tuning (SFT)} on $N=64$ samples for 25 epochs (AdamW, LR=$10^{-3}$) to test downstream compatibility.
\end{itemize}

\subsection{Main Results}
Table \ref{tab:main_results} presents the complete multi-task merging performance. We also visualize the average multi-task accuracy recovery across varying calibration budgets in Figure \ref{fig:results_budget}.

\begin{figure*}[htbp]
\centering
\includegraphics[width=0.92\textwidth]{results_plot.pdf}
\caption{Multi-task model merging accuracy recovery across different calibration dataset sizes ($N \in \{16, 64, 128\}$) under Weight Averaging (WA, left) and Task Arithmetic (TA, right). The uncalibrated merged model performance is shown as a grey dashed baseline. Our proposed S-SAC methods (\textbf{SMAC} and \textbf{C-SMAC}) achieve exceptionally high and stable accuracy across all budgets, significantly outperforming standalone spatial and frequency calibration baselines.}
\label{fig:results_budget}
\end{figure*}

\subsection{Key Findings}
\begin{enumerate}
    \item \textbf{Synergistic Dual-Domain Advantage:} S-SAC (SMAC) achieves an extraordinary performance leap, boosting uncalibrated WA average accuracy from $39.07\%$ to **$68.56\%$** ($N=64$) and TA from $40.63\%$ to **$66.65\%$** ($N=64$). This absolute improvement of over **$+26\%$** beats standalone global spatial scaling (SP-TAAC: $31.31\%$) and frequency-only alignment (FDSA: $39.64\%$) by an immense margin.
    \item \textbf{Extreme Sample Efficiency:} With only $N=16$ samples, SMAC achieves **$65.91\%$** average accuracy under WA, outperforming both SP-TAAC and FDSA using $8\times$ more data ($N=128$, achieving $38.42\%$ and $40.48\%$ respectively). This proves that joint spectral-spatial statistics are highly concentrated and can be estimated accurately with very few samples.
    \item \textbf{Head SFT Synergy:} S-SAC provides an exceptionally aligned backbone. S-SAC Standalone actually outscores Uncalibrated + Head SFT ($68.56\%$ vs. $58.37\%$) and SP-TAAC + SFT ($50.28\%$), showing that Downstream head-SFT cannot recover from bad backbone alignment, whereas S-SAC naturally resolves it.
    \item \textbf{Layer-wise vs. Channel-wise Trade-off:} Layer-wise S-SAC (SMAC) consistently outperforms Channel-wise S-SAC (C-SMAC) across all budgets (e.g., $68.56\%$ vs. $66.31\%$ under WA at $N=64$). This indicates that estimating independent 2D Fourier scaling maps for each of the $C$ channels introduces high variance under small sample regimes, whereas layer-wise magnitude scaling is highly stable and robust.
\end{enumerate}

\section{Ablation Studies and Analysis}
\label{sec:ablation}

\subsection{Sequential Calibration Path Ablation}
To verify the necessity of our dual-domain sequential path (FFT Magnitude Scaling $\rightarrow$ Spatial BatchNorm Alignment), we evaluate alternative sequential orders on WA at $N=64$:
\begin{itemize}
    \item \textbf{Spatial First, then FFT:} Under this path, we apply spatial standardization and alignment first, followed by FFT magnitude scaling. This configuration achieves $42.13\%$ average accuracy (a loss of $-26.43\%$ compared to S-SAC). This is because FFT magnitude scaling modifies the overall spatial variance, which invalidates the prior spatial normalization, demonstrating that spatial statistics alignment must always be applied \textit{after} frequency-domain magnitude adjustments.
    \item \textbf{FFT Only (FDSA):} Operating without spatial normalization achieves $39.64\%$, demonstrating that spatial alignment is absolutely vital to ground the rescaled activations.
\end{itemize}

\subsection{Fourier Phase Desynchronization Analysis}
A key question is why we preserve the merged model's phase relationships instead of aligning them with the experts' phase spectrum. We evaluated an explicit phase alignment variant (\textbf{SMAC-PR}), which rotates the merged model's Fourier phases to match the circular mean of the experts' phases.
Under WA at $N=64$, SMAC-PR catastrophically collapses, achieving exactly $10.36\%$ average accuracy (random guessing). This is because the spatial phase encodes coordinate-specific structural layout. Since different calibration images contain completely different objects at different coordinates, their phase profiles are unsynchronized. Averaging their complex phasors leads to destructive circular interference, creating a random phase mask that completely scrambles the features and destroys semantic representations. Thus, phase must remain uncalibrated.

\subsection{Sensitivity to Spectral Clamping Threshold $\gamma_{\text{max}}$}
To evaluate the impact of the maximum scaling threshold $\gamma_{\text{max}}$ on representation recovery, we perform a hyperparameter sweep over $\gamma_{\text{max}} \in \{1.0, 1.5, 2.0, 3.0, 5.0, 10.0, 20.0, 50.0\}$ on Weight Averaging with $N=64$ calibration samples. Note that setting $\gamma_{\text{max}} = 1.0$ completely clamps all spectral scaling maps to unity, thereby reducing S-SAC to a purely spatial-domain channel-wise alignment scheme. 

As shown in our empirical sweep, $\gamma_{\text{max}} = 1.0$ (spatial-only alignment) achieves $68.77\%$ average accuracy. Introducing mild spectral magnitude scaling at $\gamma_{\text{max}} = 2.0$ improves the accuracy to $69.03\%$, which further peaks at **$69.23\%$** for $\gamma_{\text{max}} = 5.0$. For larger thresholds ($\gamma_{\text{max}} \ge 10.0$), the average accuracy remains highly stable (fluctuating marginally between $69.18\%$ and $69.22\%$). This empirical stability demonstrates that S-SAC is highly robust to the exact choice of clamping factor and does not suffer from high-frequency noise amplification even under extremely relaxed clamping constraints.

\subsection{Depth-wise Calibration Ablation}
To understand where representation collapse is most severe and which network depths benefit most from dual-domain calibration, we perform a depth-wise ablation of S-SAC. We segment the 20 BatchNorm layers of ResNet-18 into three functional depth divisions:
\begin{itemize}
    \item \textbf{Early Layers Only:} Calibrating only the input convolution's BatchNorm and block-layer 1 (5 BNs total).
    \item \textbf{Middle Layers Only:} Calibrating only block-layer 2 and block-layer 3 (10 BNs total).
    \item \textbf{Late Layers Only:} Calibrating only block-layer 4 (5 BNs total).
    \item \textbf{All Layers:} Calibrating all 20 BNs (standard S-SAC / SMAC).
\end{itemize}
We run this depth ablation at $N=64$ across 3 random seeds and evaluate under both Weight Averaging (WA) and Task Arithmetic (TA) paradigms.

Under WA, calibrating only Early Layers yields an average accuracy of $55.99 \pm 0.04\%$. Calibrating only Middle Layers substantially improves recovery to $64.94 \pm 0.23\%$, while calibrating only Late Layers achieves a comparable $64.66 \pm 0.22\%$. We observe a highly consistent trend under TA, where Early, Middle, and Late calibration achieve $57.20 \pm 0.23\%$, $65.75 \pm 0.41\%$, and $62.48 \pm 0.54\%$ respectively.

These results reveal several key insights. First, calibrating early layers alone is insufficient because representation collapse compounds and amplifies as activations propagate deeper through the uncalibrated segments of the network. Second, the middle and late layers represent critical representation bottlenecks where both spectral decay and spatial distribution shifts are highly pronounced, making their alignment particularly impactful. Finally, calibrating all layers jointly achieves the highest synergistic performance ($68.85 \pm 0.47\%$ under WA; $67.39 \pm 0.33\%$ under TA). This demonstrates that spectral-spatial degradation accumulates globally throughout the model hierarchy and requires global rectification.

\subsection{Robustness to Class-Skewed Calibration Data}
In practical deployment scenarios, balanced calibration datasets are often unavailable, and the provided samples may exhibit severe distribution shifts or class imbalance. To rigorously evaluate the stability of our dual-domain framework under extreme data bias, we design a highly challenging class-skew ablation. Specifically, we restrict the $N = 64$ calibration samples for each task to contain \textit{only} classes 0--4 (the first five classes, e.g., only digits 0--4 for MNIST and early categories for FashionMNIST and CIFAR-10), completely omitting classes 5--9 during calibration. We run this severe skew experiment over 3 independent random seeds and compare S-SAC against spatial-only and frequency-only baselines.

As shown in our empirical evaluation on Weight Averaging, the uncalibrated model achieves $39.07\%$ average accuracy. Under severe class skew, the spatial-only SP-TAAC method collapses catastrophically to $34.26 \pm 0.66\%$, failing to outperform the uncalibrated baseline. This is because spatial channel-wise affine alignment attempts to force the activation distribution of skewed classes to match the balanced expert statistics, thereby distorting the representations of unseen classes. Conversely, the frequency-only FDSA baseline is less affected by semantic skew, achieving $41.48 \pm 0.67\%$ average accuracy. 

Our proposed dual-domain methods exhibit extraordinary robustness under extreme class bias. S-SAC (SMAC) achieves **$62.40 \pm 1.26\%$** average accuracy, and Channel-wise S-SAC (C-SMAC) achieves **$59.08 \pm 1.23\%$** average accuracy. Both methods outperform the spatial-only baseline by over **$+28\%$** and the frequency-only baseline by over **$+17\%$** absolute accuracy. This immense robustness demonstrates that the low-pass filter effect and spectral collapse are highly consistent, global properties of neural representation space that can be effectively and safely recovered using S-SAC even when the calibration data suffers from severe semantic skew.

\section{Conclusion, Limitations and Future Work}
In this work, we deconstructed representation collapse in model merging as an entangled spectral-spatial degradation. We proposed \textbf{Spectral-Spatial Activation Calibration (S-SAC / SMAC)}, which sequentially aligns 2D Fourier magnitude profiles and spatial channel statistics. S-SAC achieves extraordinary improvements, boosting multi-task performance by over $+29\%$ absolute accuracy on standard vision benchmarks with minimal data. S-SAC establishes a powerful new dual-domain paradigm for post-merge representation alignment.

\textbf{Limitations and Future Work.} While S-SAC demonstrates exceptional performance and data efficiency, some limitations remain. First, we evaluated S-SAC on convolutional architectures (ResNet-18) using 2D Fast Fourier Transforms. For modern transformer architectures (e.g., Vision Transformers or Large Language Models), activation sequences are 1D token sequences, and applying Fourier spectral scaling requires formulating a 1D sequence-level spectral alignment or analyzing the frequency components of self-attention matrices. Expanding S-SAC to sequence models and self-attention operations is an exciting avenue for future work. Second, our method requires running a forward pass on a small calibration set to collect expert activation statistics. Although $N=16$ or $64$ is extremely cheap, developing a zero-shot, data-free spectral-spatial calibration method based on weight statistics alone remains a valuable and open research direction.

\bibliography{submission}
\bibliographystyle{icml2026}

\end{document}
